\documentclass[12pt]{WHOIletter}
\usepackage[margin=0.5in]{geometry}  % Sets all margins to 1 inch
\WHOIname{Anthony Meza}
\position{PhD Candidate}
\officenum{21}
\email{anthony.meza@whoi.edu}
\usepackage{graphicx}

\begin{document}


\begin{letter}{
Journal of Geophysical Research (JGR) \\
Editors of {\it JGR: Oceans} \\
}
\opening{Dear Editors of {\it JGR: Oceans},} 

Please find under this cover a manuscript entitled, ``Wind Driven
Mid-Depth Pacific Cooling in a Dynamically Consistent Ocean State
Estimate,'' by Anthony Meza and Geoffrey Gebbie.  Recent comparisons
of reanalyses and sparse ship-based hydrography have highlighted
uncertainty in the sign of deep ocean temperature trends. Here, we use
the Estimating the Circulation and Climate of the Ocean (ECCO) V4r4
state estimate to design and run a suite of perturbation experiments
that quantify the dominant sources of temperature variability in the
mid-depth Pacific Ocean. Our findings indicate that wind stress
control adjustments predominantly drives cooling in ECCO V4r4 by
accelerating vertical velocities, which, in turn, causes isopycnal
heaving. \textbf{In summary, our work does not rule out the presence
  of Little Ice Age signals in the real ocean but emphasizes the need
  to consider wind-forced variability when interpreting temperature
  trends from sparse hydrographic observations.}
% \textbf{In summary, our work 
% emphasizes the need to 
% account for wind-forced variability when interpreting our sparse 
% records of deep ocean temperature trends.}

Before submitting this manuscript, we extensively reviewed and
presented portions of this work at multiple conferences and workshops
including the 2022 AGU Fall Meeting and the 2023 ECCO ECCO Annual
Meeting. The manuscript was also reviewed by a thesis committee
comprising five Physical Oceanography faculty members from MIT and
WHOI.  One reviewer and thesis committee member was Christopher
Piecuch--a member of various NASA Science Teams and an active
collaborator of the ECCO project--who provided detailed and welcomed
feedback that helped improve this manuscript.  In his view, this
manuscript ``highlights the challenges and requirements necessary to
realistically simulate deep ocean temperature trends.''

This work could fit in the scope of either the {\it JGR: Oceans} or
the {\it Journal of Advances in Modeling Earth Systems}.  Among the
two, {\it JGR: Oceans} has published more work relevant to ocean heat
content, ocean heat transport and evaluations of the ECCO state
estimate, making it a strong fit for this manuscript.  Editors with
expertise in these topics, such as Lisa Beal and Karina von
Schuckmann, further support its relevance to this journal.

% The submitted manuscript appears to be over the AMS page limit. Another internal reviewer has noted that the material is complicated due to the intricacies of the equation of state of seawater, but that the manuscript is clearly written. To accomplish this, I have streamlined the main text so that it fits within the standard ~25 page word limit, but this space is not sufficient to explain to the interested reader how to actually do the calculations. To remedy this, I have included 3 appendices that cover 5 pages to make the manuscript complete at the cost of going over the page limit. I hope the editors understand and approve the rationale behind this presentation strategy. 

% Suggestions for reviewers with expertise in this research area are: 

Thank you for considering this work for publication.
% I have used the updated LaTex template and included a Data Availability statement. 

\closing{Sincerely,}
\WHOIname\\PhD Candidate\\Woods Hole Oceanographic Institution

\end{letter}
\end{document}

% Spell check & quality check:
% detex file.tex | aspell list | sort | uniq
%
% QC:
% detex file.tex > text.txt
% diction, style, word_dupe, word_passive, word_weasel
%

%%% Local Variables:
%%% mode: LaTeX
%%% TeX-master: t
%%% End:
